% PDF 1.4 without object streams, for older viewers.
\pdfminorversion=4
\pdfobjcompresslevel=0
\documentclass[11pt]{article}
\usepackage[margin=2.3cm]{geometry}
\usepackage{amsmath,amssymb}
\usepackage{booktabs,array,longtable}
\usepackage[T1]{fontenc}
\usepackage{lmodern}
\usepackage{microtype}
\usepackage[hidelinks]{hyperref}

\renewcommand{\P}{\mathcal{P}}
\newcommand{\F}{\mathbb{F}_3}
\newcommand{\A}{\mathcal{A}}
\DeclareMathOperator{\LT}{LT}
\DeclareMathOperator{\Ext}{Ext}
\setlength{\parindent}{0pt}
\setlength{\parskip}{0.4em}

\title{A small minimal resolution over the Steenrod algebra at $p = 3$,\\ computed by Gröbner bases}
\author{Trace of the \texttt{groebner} crate (branch \texttt{groebner\_lin})}
\date{}

\begin{document}
\maketitle

\begin{abstract}
We follow, one reduction at a time, the computation of the minimal free resolution
$\dots \to C_3 \to C_2 \to C_1 \to C_0 \to \F$ of the ground field over the mod $3$ Steenrod
algebra $\A$ by the \texttt{groebner} crate, a port of Weinan Lin's Gröbner basis method
extended to odd primes. The range is homological degree $s \le 3$ and stem $t - s \le 10$, so
internal degrees $t \le 13$. Compared with $p = 2$, there are three new features: the exterior
generators $Q_k$, coefficients in $\F$, and polynomial generators of height $3$ in the
associated graded algebra. Every formula below is taken from the program's own data: the
critical pairs it formed, the elements it reduced, and the bases and differentials it
produced. Only the bookkeeping of the reduction steps was replayed outside the program, and the
replay was checked to reproduce the program's results exactly.
\end{abstract}

\section{Setting}

\paragraph{The algebra.}
We work in the Milnor basis $Q(E)\,\P(R) = Q_0^{\epsilon_0}Q_1^{\epsilon_1}\cdots
\P(r_1, r_2, \dots)$ of $\A$, writing $\P^r = \P(r)$. Here $Q_0 = \beta$ is the Bockstein,
$|Q_k| = 2\cdot 3^k - 1$, and $|\P(R)| = \sum_n 2 r_n (3^n - 1)$. Thus $|Q_0| = 1$, $|Q_1| = 5$,
$|\P^1| = 4$, and $|\P(0,1)| = 16$. Write $r_n = \sum_i d_{i,n} 3^i$ with ternary digits
$d_{i,n} \in \{0, 1, 2\}$. The digit records the power of $\xi_n^{3^i}$ in the dual basis,
of degree $2\cdot 3^i(3^n - 1)$. We record $Q(E)\,\P(R)$ by the monomial
\[ Q(E) \prod x_{in}^{d_{i,n}} \]
in the associated graded algebra of the May filtration
\[ E^0\A = E[Q_0, Q_1, \dots] \otimes \bigotimes_{i \ge 0,\, n \ge 1} \F[x_{in}]/(x_{in}^3). \]
For example $\P^2 \leftrightarrow x_{01}^2$, $\P^3 \leftrightarrow x_{11}$ and
$Q_0\P^1 \leftrightarrow Q_0x_{01}$. A monomial is given by a set of $Q_k$ and a digit
$0 \le d \le 2$ in each position. Divisibility is inclusion of the sets together with digitwise
$\le$; quotients, greatest common divisors and least common multiples are digitwise too. The
product in $E^0\A$ is non-zero exactly when no $Q_k$ repeats and no digit sum reaches $3$. Even
then it carries a coefficient: a binomial coefficient for each digit, and a sign from
reordering the $Q_k$. The positions in this range, in increasing degree, are
\[ x_{01}\ (4),\quad x_{11}\ (12),\quad x_{02}\ (16). \]

\paragraph{The monomial order.}
A term of a free module is $c\, Q(E)\P(R)\, g$ with $c \in \F^\times$ and $g$ a generator.
Terms are compared by
\begin{enumerate}
\item the generator: the generator added later (of higher degree) comes first;
\item the May weight $\sum_{k \in E}(2k+1) + \sum d_{i,n}(2n - 1)$: smaller comes first;
\item the exterior part, read as a binary number with $Q_0$ most significant: smaller comes
  first;
\item the digits, read as a number with the position of lowest degree most significant:
  smaller comes first.
\end{enumerate}
The \emph{leading term} $\LT(x)$ is the first term in this order, and elements are always
printed with it first. The leading term of a product $Q(E)\P(R)\cdot x$ is the product of
$Q(E)\,x^R$ with $\LT(x)$ in $E^0\A$, \emph{including its coefficient}, whenever that product
is non-zero. Elements of a Gröbner basis are normalized to leading coefficient $1$. Since
$p = 3$, coefficients are printed as $\pm 1$.

\paragraph{Names.}
The generator of $C_0$ is $\iota$. The generators of $C_s$, $s \ge 1$, are named after the
classes of $\Ext_\A^{s,t}(\F, \F)$ they represent:
\begin{itemize}
\item $[a_0] \in C_1$ (dual to $Q_0$) and $[h_0] \in C_1$ (dual to $\P^1$);
\item the tower $[a_0^s]$;
\item $[a_1]$ and $[b_0]$ in $C_2$;
\item $[a_0b_0]$ in $C_3$.
\end{itemize}
These are only labels. The program numbers the generators of each $C_s$ in the order they are
found.

\section{The algorithm}

The algorithm is the same as at $p = 2$. We recall it briefly, with the changes for odd primes.

\paragraph{Levels.}
Level $s$ is a Gröbner basis $b^s_0, b^s_1, \dots$ of the graph of the differential,
\[ \Gamma_s = \{ (d y, y) : y \in C_{s+1} \} \subseteq C_s \oplus C_{s+1}, \]
for the order on the first component. Each $b^s_k = (x_1 \mid x_2)$ satisfies $d x_2 = x_1$, its
leading term is that of $x_1$, and its leading coefficient is $1$. The module being resolved
enters as ``level $-1$'': $\F = \A/\A\{Q_0, \P^1, \P^3, \P^9, \dots\}$, so the relations are
$Q_0\iota = 0$ and $\P^{3^i}\iota = 0$.

\paragraph{Critical pairs and singles.}
When an element with leading term $L$ joins level $s$, the program forms
\begin{itemize}
\item a \emph{critical pair} with each earlier element $b$ whose leading term $L'$ lies on the
  same generator. Let $M = \operatorname{lcm}(L, L')$ in $E^0\A$, with cofactors $u = M/L'$ and
  $u' = M/L$. The products $u\cdot b$ and $u'\cdot b_{\text{new}}$ have leading terms
  $\lambda M$ and $\mu M$. The S-element is $\lambda\, u'\cdot b_{\text{new}} - \mu\, u\cdot b$,
  in which the leading terms cancel. The pairs are pruned by the Gebauer--Möller criteria;
\item a \emph{critical single} $m \cdot b_{\text{new}}$ for each generator $m$ of the
  annihilator of $L$ in $E^0\A$. These are $m = x_{in}^{3-d}$ for each digit $d$ of $L$ in
  position $x_{in}$, and $m = Q_k$ for each $Q_k$ in $L$. Then $m \cdot L = 0$ in $E^0\A$, so
  the product must be reduced.
\end{itemize}

\paragraph{One degree.}
In degree $t$, the program does the following:
\begin{enumerate}
\item Every level hands out its pairs and singles of degree $t$, and the relations of degree
  $t$ are taken up.
\item Each of them is reduced. The first component is reduced modulo level $s$: a term
  $c\,m$ divisible by $\LT(b)$, with quotient $u$, is cancelled by adding $-c/\lambda$ times
  $u \cdot b$, where $\lambda$ is the leading coefficient of $u \cdot b$. The second component
  is carried along. Then the second component is reduced modulo the first components of level
  $s+1$, i.e.\ modulo the boundaries known so far. If the first component has become zero, the
  second component $z \in C_{s+1}$ is a cycle.
\item Going down from the top level, the results are compared with the elements found earlier
  in the same degree, normalized, and sorted. A non-zero first component gives a \textbf{new
  basis element}. If both components are zero, the reduction was \textbf{useless}. Otherwise
  $z$ is a cycle that is not a boundary, and it becomes $d$ of a \textbf{new generator} of
  $C_{s+2}$, or the outcome is \textbf{beyond} if $s + 2 > 3$.
\end{enumerate}

\paragraph{Truncation.}
We want stems $\le 10$. A cycle found at level $s$ in degree $t$ is a generator of $C_{s+2}$ in
stem $t-s-2$, so level $s$ is only needed for $t \le s + 12$. Relations are only needed for
$t \le 11$, which leaves out $\P^3\iota$ in degree $12$ (it would give $h_1$ in stem $11$).

\paragraph{Running it.}\sloppy
The trace was produced by a small driver around the public API of the crate. It runs
\texttt{Resolution::sphere} at $p = 3$ with \texttt{SteenrodKind::Classical} and
\texttt{Bounds \{ s\_max: 3, n\_max: 10 \}}, calls \texttt{compute\_through(t)} for
$t = 1, \dots, 13$, and records every reduction through a \texttt{PairObserver}.

\section{Step by step}

\subsection*{Degree $t = 1$}
\noindent The work of this degree: the relation $Q_0\iota$. All of them are reduced independently (in parallel in the program), and the results are then sorted from the top level down.
\medskip
\paragraph{Relation $Q_0\iota = 0$.} The module relation enters as the element $(0 \mid Q_0\iota)$ of level $-1$, so its $C_0$ component is a candidate for $d$ of a new generator. It is reduced modulo the image $\operatorname{im}(d\colon C_1\to C_0)$, i.e.\ the first components of level $0$.
No term of it is divisible by a leading term of level $0$ from earlier degrees.
\textbf{Outcome: new generator.} The relation $Q_{0}\iota$ is not a boundary, so $C_{1}$ gets a generator $[a_0]$ in degree $1$ with $d[a_0] = Q_{0}\iota$.
Level $0$ gets the element $b^{0}_{0} = \bigl(Q_{0}\iota \;\big|\; [a_0]\bigr)$, with leading term $Q_{0}\,\iota$.
\medskip\noindent\emph{End of degree 1.} New basis elements: level $0$: $b^{0}_{0}$. New generators: $[a_0]\in C_{1}$.
\subsection*{Degree $t = 2$}
\noindent The work of this degree: level $0$: single. All of them are reduced independently (in parallel in the program), and the results are then sorted from the top level down.
\medskip
\paragraph{Level $0$, single $Q_{0}\cdot b^{0}_{0}$.} In $E^0\A$ we have $Q_{0}^2 = 0$, so the leading term of $Q_{0}\cdot b^{0}_{0}$ is not the product of the leading terms, and the product has to be reduced:
\[\bigl(0 \;\big|\; Q_{0}[a_0]\bigr).\]
The first component is zero, so the second component $z = Q_{0}[a_0]$ is a cycle in $C_{1}$.
No term of it is divisible by a leading term of level $1$ from earlier degrees.
The cycle $Q_{0}[a_0]$ is set aside: cycles are tested once all results of level $0$ in this degree are sorted.
\medskip
\paragraph{The cycles of level $0$.} There is one.
\textbf{Outcome: new generator.} The cycle $Q_{0}[a_0]$ is not a boundary, so $C_{2}$ gets a generator $[a_0^2]$ in degree $2$ with $d[a_0^2] = Q_{0}[a_0]$.
Level $1$ gets the element $b^{1}_{0} = \bigl(Q_{0}[a_0] \;\big|\; [a_0^2]\bigr)$, with leading term $Q_{0}\,[a_0]$.
\medskip\noindent\emph{End of degree 2.} New basis elements: level $1$: $b^{1}_{0}$. New generators: $[a_0^2]\in C_{2}$.
\subsection*{Degree $t = 3$}
\noindent The work of this degree: level $1$: single. All of them are reduced independently (in parallel in the program), and the results are then sorted from the top level down.
\medskip
\paragraph{Level $1$, single $Q_{0}\cdot b^{1}_{0}$.} In $E^0\A$ we have $Q_{0}^2 = 0$, so the leading term of $Q_{0}\cdot b^{1}_{0}$ is not the product of the leading terms, and the product has to be reduced:
\[\bigl(0 \;\big|\; Q_{0}[a_0^2]\bigr).\]
The first component is zero, so the second component $z = Q_{0}[a_0^2]$ is a cycle in $C_{2}$.
No term of it is divisible by a leading term of level $2$ from earlier degrees.
The cycle $Q_{0}[a_0^2]$ is set aside: cycles are tested once all results of level $1$ in this degree are sorted.
\medskip
\paragraph{The cycles of level $1$.} There is one.
\textbf{Outcome: new generator.} The cycle $Q_{0}[a_0^2]$ is not a boundary, so $C_{3}$ gets a generator $[a_0^3]$ in degree $3$ with $d[a_0^3] = Q_{0}[a_0^2]$.
Level $2$ gets the element $b^{2}_{0} = \bigl(Q_{0}[a_0^2] \;\big|\; [a_0^3]\bigr)$, with leading term $Q_{0}\,[a_0^2]$.
\medskip\noindent\emph{End of degree 3.} New basis elements: level $2$: $b^{2}_{0}$. New generators: $[a_0^3]\in C_{3}$.
\subsection*{Degree $t = 4$}
\noindent The work of this degree: the relation $\P^{1}\iota$; level $2$: single. All of them are reduced independently (in parallel in the program), and the results are then sorted from the top level down.
\medskip
\paragraph{Level $2$, single $Q_{0}\cdot b^{2}_{0}$.} In $E^0\A$ we have $Q_{0}^2 = 0$, so the leading term of $Q_{0}\cdot b^{2}_{0}$ is not the product of the leading terms, and the product has to be reduced:
\[\bigl(0 \;\big|\; Q_{0}[a_0^3]\bigr).\]
The first component is zero, so the second component $z = Q_{0}[a_0^3]$ is a cycle in $C_{3}$.
There is no level $3$ (we stop at $s = 3$).
The cycle $Q_{0}[a_0^3]$ is set aside: cycles are tested once all results of level $2$ in this degree are sorted.
\medskip
\paragraph{The cycles of level $2$.} There is one.
\textbf{Outcome: beyond.} The cycle $Q_{0}[a_0^3]$ would give a generator of $C_{4}$, beyond $s = 3$, so it is dropped.
\medskip
\paragraph{Relation $\P^{1}\iota = 0$.} The module relation enters as the element $(0 \mid \P^{1}\iota)$ of level $-1$, so its $C_0$ component is a candidate for $d$ of a new generator. It is reduced modulo the image $\operatorname{im}(d\colon C_1\to C_0)$, i.e.\ the first components of level $0$.
No term of it is divisible by a leading term of level $0$ from earlier degrees.
\textbf{Outcome: new generator.} The relation $\P^{1}\iota$ is not a boundary, so $C_{1}$ gets a generator $[h_0]$ in degree $4$ with $d[h_0] = \P^{1}\iota$.
Level $0$ gets the element $b^{0}_{1} = \bigl(\P^{1}\iota \;\big|\; [h_0]\bigr)$, with leading term $x_{01}\,\iota$.
\medskip\noindent\emph{End of degree 4.} New basis elements: level $0$: $b^{0}_{1}$. New generators: $[h_0]\in C_{1}$.
\subsection*{Degree $t = 5$}
\noindent The work of this degree: level $0$: pair. All of them are reduced independently (in parallel in the program), and the results are then sorted from the top level down.
\medskip
\paragraph{Level $0$, pair $(b^{0}_{0}, b^{0}_{1})$.} The leading terms $\LT(b^{0}_{0}) = Q_{0}\,\iota$ and $\LT(b^{0}_{1}) = x_{01}\,\iota$ have least common multiple $Q_{0}x_{01}\,\iota$, with cofactors $x_{01}$ and $Q_{0}$.
The leading terms of $\P^{1}\cdot b^{0}_{0}$ and $Q_{0}\cdot b^{0}_{1}$ are $1$ and $1$ times this monomial, so the S-element is $Q_{0}\cdot b^{0}_{1} - \P^{1}\cdot b^{0}_{0}$:
\[\bigl(-Q_{1}\iota \;\big|\; Q_{0}[h_0] - \P^{1}[a_0]\bigr).\]
No term of the first component is divisible by a leading term of level $0$.
Multiply by $-1$ to make the leading coefficient $1$.
\textbf{Outcome: new basis element} $b^{0}_{2} = \bigl(Q_{1}\iota \;\big|\; -Q_{0}[h_0] + \P^{1}[a_0]\bigr)$ of level $0$, with leading term $Q_{1}\,\iota$.
\medskip\noindent\emph{End of degree 5.} New basis elements: level $0$: $b^{0}_{2}$. No new generators.
\subsection*{Degree $t = 6$}
\noindent The work of this degree: level $0$: pair. All of them are reduced independently (in parallel in the program), and the results are then sorted from the top level down.
\medskip
\paragraph{Level $0$, pair $(b^{0}_{0}, b^{0}_{2})$.} The leading terms $\LT(b^{0}_{0}) = Q_{0}\,\iota$ and $\LT(b^{0}_{2}) = Q_{1}\,\iota$ have least common multiple $Q_{0}Q_{1}\,\iota$, with cofactors $Q_{1}$ and $Q_{0}$.
The leading terms of $Q_{1}\cdot b^{0}_{0}$ and $Q_{0}\cdot b^{0}_{2}$ are $-1$ and $1$ times this monomial, so the S-element is $-Q_{0}\cdot b^{0}_{2} - Q_{1}\cdot b^{0}_{0}$:
\[\bigl(0 \;\big|\; -Q_{0}\P^{1}[a_0] - Q_{1}[a_0]\bigr).\]
The first component is zero, so the second component $z = -Q_{0}\P^{1}[a_0] - Q_{1}[a_0]$ is a cycle in $C_{1}$.
Reduce it modulo $\operatorname{im}(d\colon C_{2}\to C_{1})$, the first components of level $1$:
\begin{itemize}
\item The term $-Q_{0}\P^{1}[a_0]$ has leading monomial $Q_{0}x_{01}\,[a_0]$, divisible by $\LT(b^{1}_{0})$ with quotient $x_{01}$. Adding $\P^{1}\cdot \pi_1(b^{1}_{0}) = Q_{0}\P^{1}[a_0] + Q_{1}[a_0]$ cancels it and leaves $0$.
\end{itemize}
\textbf{Outcome: useless.} Both components vanish, and the reduction contributed nothing.
\medskip\noindent\emph{End of degree 6.} No new generators.
\subsection*{Degree $t = 7$}
\noindent There are no critical pairs, singles or relations in this degree. Nothing happens.
\subsection*{Degree $t = 8$}
\noindent There are no critical pairs, singles or relations in this degree. Nothing happens.
\subsection*{Degree $t = 9$}
\noindent The work of this degree: level $0$: pair. All of them are reduced independently (in parallel in the program), and the results are then sorted from the top level down.
\medskip
\paragraph{Level $0$, pair $(b^{0}_{1}, b^{0}_{2})$.} The leading terms $\LT(b^{0}_{1}) = x_{01}\,\iota$ and $\LT(b^{0}_{2}) = Q_{1}\,\iota$ have least common multiple $Q_{1}x_{01}\,\iota$, with cofactors $Q_{1}$ and $x_{01}$.
The leading terms of $Q_{1}\cdot b^{0}_{1}$ and $\P^{1}\cdot b^{0}_{2}$ are $1$ and $1$ times this monomial, so the S-element is $\P^{1}\cdot b^{0}_{2} - Q_{1}\cdot b^{0}_{1}$:
\[\bigl(0 \;\big|\; -Q_{0}\P^{1}[h_0] + Q_{1}[h_0] - \P^{2}[a_0]\bigr).\]
The first component is zero, so the second component $z = -Q_{0}\P^{1}[h_0] + Q_{1}[h_0] - \P^{2}[a_0]$ is a cycle in $C_{1}$.
No term of it is divisible by a leading term of level $1$ from earlier degrees.
The cycle $-Q_{0}\P^{1}[h_0] + Q_{1}[h_0] - \P^{2}[a_0]$ is set aside: cycles are tested once all results of level $0$ in this degree are sorted.
\medskip
\paragraph{The cycles of level $0$.} There is one.
\textbf{Outcome: new generator.} The cycle $Q_{0}\P^{1}[h_0] - Q_{1}[h_0] + \P^{2}[a_0]$, multiplied by $-1$ to make the leading coefficient $1$, is not a boundary, so $C_{2}$ gets a generator $[a_1]$ in degree $9$ with $d[a_1] = Q_{0}\P^{1}[h_0] - Q_{1}[h_0] + \P^{2}[a_0]$.
Level $1$ gets the element $b^{1}_{1} = \bigl(Q_{0}\P^{1}[h_0] - Q_{1}[h_0] + \P^{2}[a_0] \;\big|\; [a_1]\bigr)$, with leading term $Q_{0}x_{01}\,[h_0]$.
\medskip\noindent\emph{End of degree 9.} New basis elements: level $1$: $b^{1}_{1}$. New generators: $[a_1]\in C_{2}$.
\subsection*{Degree $t = 10$}
\noindent The work of this degree: level $0$: single; level $1$: single. All of them are reduced independently (in parallel in the program), and the results are then sorted from the top level down.
\medskip
\paragraph{Level $1$, single $Q_{0}\cdot b^{1}_{1}$.} In $E^0\A$ we have $Q_{0}^2 = 0$, so the leading term of $Q_{0}\cdot b^{1}_{1}$ is not the product of the leading terms, and the product has to be reduced:
\[\bigl(-Q_{0}Q_{1}[h_0] + Q_{0}\P^{2}[a_0] \;\big|\; Q_{0}[a_1]\bigr).\]
Reduce the first component modulo level $1$:
\begin{itemize}
\item The term $Q_{0}\P^{2}[a_0]$ has leading monomial $Q_{0}x_{01}^{2}\,[a_0]$, divisible by $\LT(b^{1}_{0})$ with quotient $x_{01}^{2}$. Adding $-\P^{2}\cdot b^{1}_{0} = \bigl(-Q_{0}\P^{2}[a_0] - Q_{1}\P^{1}[a_0] \;\big|\; -\P^{2}[a_0^2]\bigr)$ cancels it and leaves $\bigl(-Q_{0}Q_{1}[h_0] - Q_{1}\P^{1}[a_0] \;\big|\; Q_{0}[a_1] - \P^{2}[a_0^2]\bigr)$.
\end{itemize}
Multiply by $-1$ to make the leading coefficient $1$.
\textbf{Outcome: new basis element} $b^{1}_{2} = \bigl(Q_{0}Q_{1}[h_0] + Q_{1}\P^{1}[a_0] \;\big|\; -Q_{0}[a_1] + \P^{2}[a_0^2]\bigr)$ of level $1$, with leading term $Q_{0}Q_{1}\,[h_0]$.
\medskip
\paragraph{Level $0$, single $Q_{1}\cdot b^{0}_{2}$.} In $E^0\A$ we have $Q_{1}^2 = 0$, so the leading term of $Q_{1}\cdot b^{0}_{2}$ is not the product of the leading terms, and the product has to be reduced:
\[\bigl(0 \;\big|\; Q_{0}Q_{1}[h_0] + Q_{1}\P^{1}[a_0]\bigr).\]
The first component is zero, so the second component $z = Q_{0}Q_{1}[h_0] + Q_{1}\P^{1}[a_0]$ is a cycle in $C_{1}$.
No term of it is divisible by a leading term of level $1$ from earlier degrees.
The second component contains the leading term of $b^{1}_{2}$, found in this degree; subtracting its first component leaves $0$.
\textbf{Outcome: useless.} The cycle is a multiple of the first component of $b^{1}_{2}$, found one level up earlier in this degree, so it is the boundary $d\bigl(-Q_{0}[a_1] + \P^{2}[a_0^2]\bigr)$. Both components now vanish, and the reduction contributed nothing.
\medskip\noindent\emph{End of degree 10.} New basis elements: level $1$: $b^{1}_{2}$. No new generators.
\subsection*{Degree $t = 11$}
\noindent The work of this degree: level $1$: single. All of them are reduced independently (in parallel in the program), and the results are then sorted from the top level down.
\medskip
\paragraph{Level $1$, single $Q_{0}\cdot b^{1}_{2}$.} In $E^0\A$ we have $Q_{0}^2 = 0$, so the leading term of $Q_{0}\cdot b^{1}_{2}$ is not the product of the leading terms, and the product has to be reduced:
\[\bigl(Q_{0}Q_{1}\P^{1}[a_0] \;\big|\; Q_{0}\P^{2}[a_0^2]\bigr).\]
Reduce the first component modulo level $1$:
\begin{itemize}
\item The term $Q_{0}Q_{1}\P^{1}[a_0]$ has leading monomial $Q_{0}Q_{1}x_{01}\,[a_0]$, divisible by $\LT(b^{1}_{0})$ with quotient $Q_{1}x_{01}$. Adding $Q_{1}\P^{1}\cdot b^{1}_{0} = \bigl(-Q_{0}Q_{1}\P^{1}[a_0] \;\big|\; Q_{1}\P^{1}[a_0^2]\bigr)$ cancels it and leaves $\bigl(0 \;\big|\; Q_{0}\P^{2}[a_0^2] + Q_{1}\P^{1}[a_0^2]\bigr)$.
\end{itemize}
The first component is zero, so the second component $z = Q_{0}\P^{2}[a_0^2] + Q_{1}\P^{1}[a_0^2]$ is a cycle in $C_{2}$.
Reduce it modulo $\operatorname{im}(d\colon C_{3}\to C_{2})$, the first components of level $2$:
\begin{itemize}
\item The term $Q_{0}\P^{2}[a_0^2]$ has leading monomial $Q_{0}x_{01}^{2}\,[a_0^2]$, divisible by $\LT(b^{2}_{0})$ with quotient $x_{01}^{2}$. Adding $-\P^{2}\cdot \pi_1(b^{2}_{0}) = -Q_{0}\P^{2}[a_0^2] - Q_{1}\P^{1}[a_0^2]$ cancels it and leaves $0$.
\end{itemize}
\textbf{Outcome: useless.} Both components vanish, and the reduction contributed nothing.
\medskip\noindent\emph{End of degree 11.} No new generators.
\subsection*{Degree $t = 12$}
\noindent Relations of degree $12$ would give generators of $C_1$ in stem $11 > 10$ and are no longer taken up.
\noindent The work of this degree: level $0$: single. All of them are reduced independently (in parallel in the program), and the results are then sorted from the top level down.
\medskip
\paragraph{Level $0$, single $x_{01}^{2}\cdot b^{0}_{1}$.} In $E^0\A$ we have $x_{01}^{2}\cdot x_{01} = 0$ (as $x_{in}^{3} = 0$), so the leading term of $\P^{2}\cdot b^{0}_{1}$ is not the product of the leading terms, and the product has to be reduced:
\[\bigl(0 \;\big|\; \P^{2}[h_0]\bigr).\]
The first component is zero, so the second component $z = \P^{2}[h_0]$ is a cycle in $C_{1}$.
No term of it is divisible by a leading term of level $1$ from earlier degrees.
The cycle $\P^{2}[h_0]$ is set aside: cycles are tested once all results of level $0$ in this degree are sorted.
\medskip
\paragraph{The cycles of level $0$.} There is one.
\textbf{Outcome: new generator.} The cycle $\P^{2}[h_0]$ is not a boundary, so $C_{2}$ gets a generator $[b_0]$ in degree $12$ with $d[b_0] = \P^{2}[h_0]$.
Level $1$ gets the element $b^{1}_{3} = \bigl(\P^{2}[h_0] \;\big|\; [b_0]\bigr)$, with leading term $x_{01}^{2}\,[h_0]$.
\medskip\noindent\emph{End of degree 12.} New basis elements: level $1$: $b^{1}_{3}$. New generators: $[b_0]\in C_{2}$.
\subsection*{Degree $t = 13$}
\noindent Level $0$ is no longer needed: a cycle found at level $s$ in this degree would be a generator of $C_{s+2}$ in stem $13-s-2 > 10$, so its pairs are not formed.
\noindent The work of this degree: level $1$: pair. All of them are reduced independently (in parallel in the program), and the results are then sorted from the top level down.
\medskip
\paragraph{Level $1$, pair $(b^{1}_{1}, b^{1}_{3})$.} The leading terms $\LT(b^{1}_{1}) = Q_{0}x_{01}\,[h_0]$ and $\LT(b^{1}_{3}) = x_{01}^{2}\,[h_0]$ have least common multiple $Q_{0}x_{01}^{2}\,[h_0]$, with cofactors $x_{01}$ and $Q_{0}$.
The leading terms of $\P^{1}\cdot b^{1}_{1}$ and $Q_{0}\cdot b^{1}_{3}$ are $-1$ and $1$ times this monomial, so the S-element is $-Q_{0}\cdot b^{1}_{3} - \P^{1}\cdot b^{1}_{1}$:
\[\bigl(0 \;\big|\; -Q_{0}[b_0] - \P^{1}[a_1]\bigr).\]
The first component is zero, so the second component $z = -Q_{0}[b_0] - \P^{1}[a_1]$ is a cycle in $C_{2}$.
No term of it is divisible by a leading term of level $2$ from earlier degrees.
The cycle $-Q_{0}[b_0] - \P^{1}[a_1]$ is set aside: cycles are tested once all results of level $1$ in this degree are sorted.
\medskip
\paragraph{The cycles of level $1$.} There is one.
\textbf{Outcome: new generator.} The cycle $Q_{0}[b_0] + \P^{1}[a_1]$, multiplied by $-1$ to make the leading coefficient $1$, is not a boundary, so $C_{3}$ gets a generator $[a_0b_0]$ in degree $13$ with $d[a_0b_0] = Q_{0}[b_0] + \P^{1}[a_1]$.
Level $2$ gets the element $b^{2}_{1} = \bigl(Q_{0}[b_0] + \P^{1}[a_1] \;\big|\; [a_0b_0]\bigr)$, with leading term $Q_{0}\,[b_0]$.
\medskip\noindent\emph{End of degree 13.} New basis elements: level $2$: $b^{2}_{1}$. New generators: $[a_0b_0]\in C_{3}$.

\section{The result}

\paragraph{Generators and differentials.}
\begin{center}
\begin{tabular}{lcccl}
\toprule
generator & $s$ & $t$ & $t-s$ & differential \\
\midrule
$\iota$ & 0 & 0 & 0 & --- \\
$[a_0]$ & 1 & 1 & 0 & $Q_{0}\iota$ \\
$[a_0^2]$ & 2 & 2 & 0 & $Q_{0}[a_0]$ \\
$[a_0^3]$ & 3 & 3 & 0 & $Q_{0}[a_0^2]$ \\
$[h_0]$ & 1 & 4 & 3 & $\P^{1}\iota$ \\
$[a_1]$ & 2 & 9 & 7 & $Q_{0}\P^{1}[h_0] - Q_{1}[h_0] + \P^{2}[a_0]$ \\
$[b_0]$ & 2 & 12 & 10 & $\P^{2}[h_0]$ \\
$[a_0b_0]$ & 3 & 13 & 10 & $Q_{0}[b_0] + \P^{1}[a_1]$ \\
\bottomrule
\end{tabular}
\end{center}
Since the resolution is minimal, the generators are a basis of $\Ext_\A^{s,t}(\F,\F)$ in this
range:
\begin{center}
\begin{tabular}{c|ccccccccccc}
$s \backslash t-s$ & 0 & 1 & 2 & 3 & 4 & 5 & 6 & 7 & 8 & 9 & 10 \\
\midrule
$3$ & $a_0^3$ &  &  &  &  &  &  &  &  &  & $a_0b_0$ \\
$2$ & $a_0^2$ &  &  &  &  &  &  & $a_1$ &  &  & $b_0$ \\
$1$ & $a_0$ &  &  & $h_0$ &  &  &  &  &  &  &  \\
$0$ & $1$ &  &  &  &  &  &  &  &  &  &  \\
\end{tabular}
\end{center}
This is the familiar picture at $p = 3$. There is the $a_0$-tower in stem $0$, $h_0$ in stem
$3$ (detecting $\alpha_1$), $a_1$ in stem $7$ (detecting $\alpha_2$), and $b_0$ in stem $10$
(detecting $\beta_1$). In a minimal resolution, the coefficient of $Q_0$ (respectively $\P^1$)
on a generator $g$ in $d$ of a generator $g'$ is the coefficient of $g'$ in the Yoneda product
$a_0 g$ (respectively $h_0 g$). So $d[a_0b_0] = Q_0[b_0] + \P^1[a_1]$ shows that the class in
stem $10$, $s = 3$, is $a_0 b_0 = h_0 a_1$, up to sign. The cycle beyond $s = 3$ found in
degree $4$ is $a_0^4$.

\paragraph{The Gröbner bases.}
\begin{center}
\small
\begin{longtable}{lclll}
\toprule
element & $t$ & leading term & first component $x_1$ & second component $x_2$ \\
\midrule
$b^{0}_{0}$ & 1 & $Q_{0}\,\iota$ & $Q_{0}\iota$ & $[a_0]$ \\
$b^{0}_{1}$ & 4 & $x_{01}\,\iota$ & $\P^{1}\iota$ & $[h_0]$ \\
$b^{0}_{2}$ & 5 & $Q_{1}\,\iota$ & $Q_{1}\iota$ & $-Q_{0}[h_0] + \P^{1}[a_0]$ \\
$b^{1}_{0}$ & 2 & $Q_{0}\,[a_0]$ & $Q_{0}[a_0]$ & $[a_0^2]$ \\
$b^{1}_{1}$ & 9 & $Q_{0}x_{01}\,[h_0]$ & $Q_{0}\P^{1}[h_0] - Q_{1}[h_0] + \P^{2}[a_0]$ & $[a_1]$ \\
$b^{1}_{2}$ & 10 & $Q_{0}Q_{1}\,[h_0]$ & $Q_{0}Q_{1}[h_0] + Q_{1}\P^{1}[a_0]$ & $-Q_{0}[a_1] + \P^{2}[a_0^2]$ \\
$b^{1}_{3}$ & 12 & $x_{01}^{2}\,[h_0]$ & $\P^{2}[h_0]$ & $[b_0]$ \\
$b^{2}_{0}$ & 3 & $Q_{0}\,[a_0^2]$ & $Q_{0}[a_0^2]$ & $[a_0^3]$ \\
$b^{2}_{1}$ & 13 & $Q_{0}\,[b_0]$ & $Q_{0}[b_0] + \P^{1}[a_1]$ & $[a_0b_0]$ \\
\bottomrule
\end{longtable}
\end{center}

\paragraph{Counters.}
The statistics the program keeps per level ($s = -1$ are the module relations):
\begin{center}
\footnotesize\setlength{\tabcolsep}{3.5pt}
\begin{tabular}{rrrrrrrrrrrr}
\toprule
$s$ & formed & divisible & chain & singles & reduced & useless & elements & boundary & gens & beyond & steps \\
\midrule
-1 & 0 & 0 & 0 & 0 & 2 & 0 & 0 & 0 & 2 & 0 & 0 \\
0 & 3 & 0 & 0 & 3 & 6 & 2 & 1 & 0 & 3 & 0 & 1 \\
1 & 1 & 0 & 0 & 3 & 4 & 1 & 1 & 0 & 2 & 0 & 3 \\
2 & 0 & 0 & 0 & 1 & 1 & 0 & 0 & 0 & 0 & 1 & 0 \\
all & 4 & 0 & 0 & 7 & 13 & 3 & 2 & 0 & 7 & 1 & 4 \\
\bottomrule
\end{tabular}
\end{center}
Here ``formed'' counts critical pairs, ``divisible'' and ``chain'' those removed by the
Gebauer--Möller criteria (none in this small range), and ``steps'' the reduction steps. Of the
13 reductions, 3 were useless and one was beyond the range. The other 9 each produced a basis
element or a generator.

\end{document}
